\section{Theoretical properties}
\label{sec:theory}

This section collects the formal properties of the test. We establish four
things: that the Monte Carlo test is finite-sample valid whenever the realized
schedule is exchangeable with its $\mathbb{Q}_\theta$-resamples; that, under a
named high-level central-limit condition, it is consistent, with a minimum
detectable edge of order $N^{-1/2}$ in the trade count; that, under the same
condition, it attains the local power of the one-sided Gaussian test against
placement alternatives while carrying \emph{zero} noncentrality against the pure
exposure direction; and that, because the conditioning measure is a modeling
choice, the timing estimand is reported only as a sensitivity range over a
declared admissible menu of measures, for whose boundary---measure-invariant
skill---a valid intersection--union test exists. Throughout, $T(S)$ is the
cumulative-return statistic induced by a schedule $S$, $S_0$ is the realized
schedule, and $\mathcal C=\sigma(\mathcal S,\{P_t\})$ is the conditioning
$\sigma$-field of Section~\ref{sec:conditioning-estimand}. We are explicit about
which results are exact, which hold only under the stated regularity condition,
and which are definitional.

\subsection{Finite-sample validity}
\label{sec:fs-validity}

The validity of the procedure rests on no model for returns, but on the
analyst's choice of reference law $\mathbb{Q}_\theta$: it requires only that the
realized statistic be exchangeable with its resampled counterparts under that
chosen law. We restate the result at the level of the statistic, so that the
section is self-contained and the conditions are visible.

\begin{proposition}[Finite-sample validity and large-$M$ consistency of the
Monte Carlo $p$-value]
\label{prop:q-validity}
Fix a measure $\mathbb{Q}_\theta$ in the family of Definition~\ref{def:h0}, and
let $S_1,\dots,S_M$ be i.i.d.\ draws from $\mathbb{Q}_\theta$ given $\mathcal C$.
Define
\[
\hat p_M
=\frac{1+\sum_{m=1}^{M}\mathbf 1\{T(S_m)\ge T(S_0)\}}{M+1}.
\]
\emph{(i) Validity.} Suppose $H_0^\theta$ holds in the statistic-level form that
$T(S_0)$ is exchangeable with $\{T(S_m)\}_{m=1}^{M}$ given $\mathcal C$; this
holds in particular if $S_0\stackrel{d}{=}\mathbb{Q}_\theta$ given $\mathcal C$.
Then $\hat p_M$ is super-uniform: for every $\alpha$ on the grid
$\{1/(M+1),\dots,1\}$, $\Pr(\hat p_M\le\alpha\mid\mathcal C)\le\alpha$, with
equality absent ties and strict conservatism when ties are counted into the
right tail.
\emph{(ii) Consistency in $M$.} For \emph{any} fixed $S_0$, whether or not
$H_0^\theta$ holds, $\hat p_M\to q_\theta:=\mathbb{Q}_\theta\big(T(S)\ge
T(S_0)\mid\mathcal C\big)$ almost surely as $M\to\infty$.
\end{proposition}

\begin{proof}
(i) Exchangeability of $T(S_0)$ with $\{T(S_m)\}_{m=1}^{M}$ given $\mathcal C$
makes the rank of $T(S_0)$ among $\{T(S_m)\}_{m=0}^{M}$ uniform on
$\{1,\dots,M+1\}$ absent ties; this holds whenever $S_0\stackrel{d}{=}
\mathbb{Q}_\theta$ given $\mathcal C$, because then $S_0,S_1,\dots,S_M$ are
themselves exchangeable and $T$ is a fixed measurable map. Dividing the
upper-tail rank by $M+1$ and counting ties into the right tail gives
$\Pr(\hat p_M\le\alpha\mid\mathcal C)\le\alpha$. (ii) For fixed $S_0$ the
indicators $\mathbf 1\{T(S_m)\ge T(S_0)\}$, $m\ge1$, are i.i.d.\
Bernoulli$(q_\theta)$ given $\mathcal C$; the strong law of large numbers yields
the limit, which does not invoke $H_0^\theta$.
\end{proof}

Part~(i) is the finite-sample validity of a Monte Carlo randomization test:
drawing $M$ schedules rather than enumerating $\mathbb{Q}_\theta$ preserves exact
level-$\alpha$ validity \cite{Dwass1957,Hope1968}; the modern group-theoretic
account is \cite{HemerikGoeman2018}, and the $+1$ in numerator and denominator is
what makes $\hat p_M$ valid at finite $M$---the naive $\sum\mathbf
1\{\cdot\}/M$ is downward biased and can be invalidly zero
\cite{PhipsonSmyth2010}. Because the null is generated by a constrained,
dependence-preserving resampling of one realized path rather than by i.i.d.\
draws, the relevant guarantee is that of \cite{BesagClifford1989}, which is exact
for precisely this conditional-on-the-data construction. We stress two limits on
the word ``exact.'' First, exactness is \emph{conditional on the reference law}:
$\mathbb{Q}_\theta$ is an analyst-chosen modeling object, not a data-determined
one, and the guarantee transfers to the realized data only insofar as the
realized schedule is genuinely exchangeable with its $\mathbb{Q}_\theta$-resamples
under $H_0^\theta$. Second, exactness holds on the attainable $p$-grid only.
When $N$ is small the gap-permutation orbit is small, the statistic $T$ is
atomic, the achievable $p$-grid is coarse, and equality in the super-uniformity
bound fails; for such strategies---e.g.\ Squeeze Breakout with $N=4$, whose
internal-gap multiset admits only $3!$ permutations crossed with the leading-gap
draw---the procedure is exploratory, not inferential, and near-threshold
$p$-values should be read accordingly. The dependence of the rejection
probability on $M$ is honestly nonmonotone: on each plateau where
$k_\alpha=\lfloor\alpha(M+1)\rfloor-1$ is constant the rejection probability
declines, jumping upward only when $\alpha(M+1)$ crosses an integer, so the
finite-$M$ power traces a sawtooth around its $M\to\infty$ limit
\cite{Hope1968}; at $M=5{,}000$ and $\alpha=0.05$ this budget is far from binding.

\subsection{A central-limit condition for the standardized statistic}
\label{sec:clt-assumption}

The consistency, minimum-detectable-edge, and local-power results below are
asymptotic, and they require a limit theorem for the standardized statistic. We
do not have an i.i.d.\ structure to invoke: under $\mathbb{Q}_\theta$ the trade
contributions are a gap-permutation of one fixed, autocorrelated return path, so
the summands are neither independent nor identically distributed, and the
placements are further coupled by the non-overlap constraint and the shared path.
We therefore state the needed Gaussian approximation as an explicit, named
high-level condition and prove the propositions \emph{conditional} on it, rather
than assert it as a by-product.

\begin{assumption}[Combinatorial / finite-population CLT condition]
\label{ass:clt}
Write $\log(1+CR)=\sum_{j=1}^{N}\xi_j$, where under $\mathbb{Q}_\theta$ the $\xi_j$
are the $N$ trade log-contributions evaluated at a placement of the fixed
duration/gap template drawn from $\mathbb{Q}_\theta$ on the fixed path. Let
$\mu_{\mathbb Q},\sigma_{\mathbb Q}$ denote the conditional $\mathbb{Q}_\theta$
mean and standard deviation of $T(S)=\log(1+CR)$ and set
$Z_N=(T(S_0)-\mu_{\mathbb Q})/\sigma_{\mathbb Q}$. We assume that under
$H_0^\theta$,
\[
Z_N \;\to_d\; N(0,1) \qquad (N\to\infty),
\]
and that this convergence holds under the following two regularity requirements:
\emph{(a)} a no-dominant-term, Lindeberg-type condition
$\max_{j}\mathrm{Var}(\xi_j)/\sum_{j}\mathrm{Var}(\xi_j)\to0$, so that no
single trade carries a non-negligible share of the return budget; and
\emph{(b)} weak ($\alpha$-mixing) dependence of the realized path, so that the
autocorrelation entering the permutation variance decays fast enough for the
finite-population variance to be the correct standardization.
\end{assumption}

This is the appropriate device. Conditional on $\mathcal C$, a
$\mathbb{Q}_\theta$-draw is a random permutation applied to a fixed array of trade
contributions, and the limiting normality of such standardized permutation
statistics is exactly the content of Hoeffding's combinatorial central limit
theorem \cite{Hoeffding1951}, whose proof turns on a no-dominant-term condition of
the Lindeberg type in part~(a). The exchangeability that licenses the conditional
construction is the classical exchangeability theory \cite{Aldous1985}, and the
finite-population permutation framework is that of \cite{LehmannRomano2005}; the
mixing requirement~(b) is what lets the fixed path's autocorrelation be absorbed
into a finite-population variance rather than break the approximation. The
condition is falsifiable and can fail: when one or a few trades dominate the
return budget, (a) is violated, the normal approximation degrades, and the
standardized statistic is not pivotal. We therefore treat $\mathrm{RCSI}_z$ as a
descriptive, approximately pivotal effect size only; the test's \emph{validity}
rests on the finite-sample exchangeability of
Proposition~\ref{prop:q-validity}(i), not on Assumption~\ref{ass:clt}, which is
invoked solely for the asymptotic statements that follow. The adequacy of the
approximation is an empirical matter, confirmed by the Kolmogorov--Smirnov
uniformity tests and the controlled-strength power calibration of
Section~\ref{sec:simulation-validation}.

\subsection{Consistency and the minimum detectable edge}
\label{sec:mde}

Validity is the floor; the test must also detect a genuine edge as evidence
accumulates. Evidence here accumulates in the trade count $N$, not in the length
of the price path, because $\mathbb{Q}_\theta$ rearranges a fixed set of $N$
trade contributions on a fixed path.

\begin{proposition}[Consistency and the $N^{-1/2}$ detection rate]
\label{prop:consistency}
Suppose Assumption~\ref{ass:clt} holds. Consider a per-trade location-shift
alternative in which the realized placement confers a common standardized entry
edge $\zeta=\mu_1/\sigma_1>0$ on each of the $N$ trade contributions to
$\log(1+CR)$, relative to a $\mathbb{Q}_\theta$-mean of zero, with per-trade
standard deviation $\sigma_1$ and $\sigma_{\mathbb Q}^2=\sum_j\mathrm{Var}(\xi_j)$.
Then $Z_N=\delta_N+O_p(1)$ with deterministic mean displacement $\delta_N=
N\mu_1/\sigma_{\mathbb Q}$, and under the no-dominant-term regime where
$\sigma_{\mathbb Q}\asymp\sqrt N\,\sigma_1$ this gives $\delta_N\asymp
\sqrt N\,\zeta$. Hence the realized $\mathbb{Q}_\theta$-percentile tends to one and
the level-$\alpha$ test is consistent as $N\to\infty$; equivalently, the minimum
edge detectable at fixed power shrinks as $\zeta_{\min}\asymp N^{-1/2}$.
\end{proposition}

\begin{proof}
Under Assumption~\ref{ass:clt}, $Z_N=(T(S_0)-\mu_{\mathbb Q})/\sigma_{\mathbb
Q}\to_d N(0,1)$ when $H_0^\theta$ holds. The location-shift alternative adds a
\emph{deterministic} mean to each contribution: $E[\log(1+CR)]$ moves from
$\mu_{\mathbb Q}$ to $\mu_{\mathbb Q}+N\mu_1$, while the centering and scaling
$(\mu_{\mathbb Q},\sigma_{\mathbb Q})$ are fixed functionals of $\mathbb{Q}_\theta$
and $\mathcal C$. Subtracting the null mean and dividing by $\sigma_{\mathbb Q}$
therefore displaces the standardized statistic by the deterministic amount
$\delta_N=N\mu_1/\sigma_{\mathbb Q}$, leaving an $O_p(1)$ fluctuation governed by
the same limit law; this is the derivation, not an assumption, of the location
shift. Writing $\delta_N=\sqrt N\,(\mu_1/\sigma_1)\cdot(\sqrt N\,\sigma_1/
\sigma_{\mathbb Q})$ and using $\sigma_{\mathbb Q}\asymp\sqrt N\,\sigma_1$
(the no-dominant-term regime of Assumption~\ref{ass:clt}(a), under which the
variance accumulates linearly in $N$) gives $\delta_N\asymp\sqrt N\,\zeta\to
\infty$. The upper-tail probability $q_\theta=\Pr_{\mathbb Q}(Z\ge Z_N)$ then
tends to zero, and Proposition~\ref{prop:q-validity}(ii) makes $\hat p_M$ track
it. Fixing the noncentrality $\sqrt N\,\zeta$ at the value delivering a target
power inverts to $\zeta_{\min}\asymp N^{-1/2}$.
\end{proof}

Two consequences match the empirical record. Low-frequency strategies are weakly
identified: Squeeze Breakout's $N=4$ admits only a large $\zeta_{\min}$, which is
why its near-threshold $p$-value is read as exploratory rather than as evidence of
absence---and, as noted in Section~\ref{sec:fs-validity}, at $N=4$ the
asymptotics of Assumption~\ref{ass:clt} are not in force at all, so the reading
is doubly cautious. And because power is governed by $N$, a non-rejection at
large $N$ is informative while a non-rejection at small $N$ is not.

\subsection{Local power and the exposure direction}
\label{sec:local-power}

The $N^{-1/2}$ boundary of Proposition~\ref{prop:consistency} is sharp: at exactly
that rate the test has nondegenerate, computable power.

\begin{proposition}[Local power against placement alternatives]
\label{prop:local-power}
Suppose Assumption~\ref{ass:clt} holds, and consider the Pitman sequence
$\zeta_N=h/\sqrt N$ for fixed $h\ge0$. Then the mean displacement of
Proposition~\ref{prop:consistency} is $\delta_N=\sqrt N\,\zeta_N+o(1)=h+o(1)$, a
bounded shift, so the sequence of alternatives is contiguous to the null and
$Z_N\to_d N(h,1)$. The one-sided level-$\alpha$ test then has asymptotic power
\[
\beta(h)=\Phi\big(h-z_{1-\alpha}\big),
\]
with $z_{1-\alpha}$ the standard normal quantile. Hence $\beta(0)=\alpha$,
$\beta(z_{1-\alpha})=\tfrac12$, and at $\alpha=0.05$ power reaches $0.80$ at
$h=z_{0.95}+z_{0.80}\approx2.49$.
\end{proposition}

\begin{proof}
By Proposition~\ref{prop:consistency} the deterministic displacement is
$\delta_N=N\mu_1/\sigma_{\mathbb Q}=\sqrt N\,\zeta_N\cdot(\sqrt N\,\sigma_1/
\sigma_{\mathbb Q})\to h$ along $\zeta_N=h/\sqrt N$, because the second factor
tends to $1$ under the no-dominant-term scaling. A bounded mean shift added to a
statistic that is asymptotically $N(0,1)$ under the null yields a contiguous
sequence (it does not drive the likelihood-ratio to a degenerate limit), and the
limit law shifts to $N(h,1)$; contiguity here is derived from the boundedness of
$\delta_N$, not assumed. Therefore $\Pr(Z_N\ge z_{1-\alpha})\to1-\Phi(z_{1-\alpha}
-h)=\Phi(h-z_{1-\alpha})$. The three values follow by substitution;
$\Phi(z_{0.80})=0.80$ gives the stated $h\approx2.49$.
\end{proof}

We now make precise the sense in which the test is orthogonal to exposure. The
key is an exact, construction-level decomposition of the log return.

\begin{proposition}[Orthogonality to exposure; shared leading noncentrality
against placement]
\label{prop:orthogonality}
Let SP denote the structure-preserving test and let RC/SPA denote the
single-strategy specialization of White's Reality Check \cite{white2000} and
Hansen's SPA test \cite{hansen2005}---studentized bootstrap tests of mean trade
return. Decompose
\[
\log(1+CR)=A(\text{path},\mathcal S)+B(\text{placement}),
\]
where the \emph{exposure} component $A$ is the return any structurally identical
schedule earns from holding the duration/gap template on the path, and the
\emph{placement} component $B$ is the excess of the realized schedule over its
$\mathbb{Q}_\theta$-resamples. By construction every $\mathbb{Q}_\theta$-resample
reprices the same trade template on the same path with the same geometry, so $A$
is \emph{invariant} across resamples: it contributes zero to the conditional
$\mathbb{Q}_\theta$-variance and zero to the noncentrality of the centered
statistic $Z_N$; only $B$ varies. Consequently, under Assumption~\ref{ass:clt}:
\emph{(i)} against a pure-exposure alternative---in which $A$ is shifted upward
while $B$ equals its $\mathbb{Q}_\theta$-mean---SP has noncentrality exactly $0$,
whereas a return-mean test loads on $E[A]+E[B]>0$ and so has strictly positive
noncentrality and rejects; \emph{(ii)} against a placement alternative under a
shared location model, SP and RC/SPA share the same leading noncentrality $h$, so
neither test dominates the other locally.
\end{proposition}

\begin{proof}
The decomposition is exact and, crucially, $A$ is $\mathcal C$-measurable up to
the placement: every resample $S_m$ holds fixed the path, the duration profile,
the gap multiset, the weights, and the directions, and reprices the identical
template, so $A(S_m)=A(S_0)$ for all $m$. Hence $A$ is constant on the
$\mathbb{Q}_\theta$-orbit given $\mathcal C$: $\mathrm{Var}_{\mathbb Q}(A\mid
\mathcal C)=0$, and the standardized statistic $Z_N=(A+B-\mu_{\mathbb Q})/
\sigma_{\mathbb Q}$ depends on the realization only through $B$, because the
constant $A$ cancels against its appearance in $\mu_{\mathbb Q}$. (i) A
pure-exposure alternative raises $E[A]$ while leaving $B$ distributed as its
$\mathbb{Q}_\theta$-law, so $B-E_{\mathbb Q}[B]$ is a central draw and the
noncentrality of $Z_N$ is zero; the profitable-but-untimed schedule is, by
construction, a typical draw under $\mathbb{Q}_\theta$. RC/SPA studentize the
realized mean trade return against a zero benchmark, so their noncentrality is the
total mean return $E[A]+E[B]$, which is strictly positive when exposure alone is
profitable. (ii) Under a shared location model the placement alternative
displaces only $B$, by the location shift of Proposition~\ref{prop:local-power};
both tests then carry the same leading noncentrality $h$ in $B$, so to first order
their local powers coincide and neither dominates. We claim only this shared
leading noncentrality and deliberately do \emph{not} assert asymptotic relative
efficiency equal to one: equality of the full local power functions would require
matching the tests' variance functionals against placement, which holds only
under the shared-location model and which we do not establish in general.
\end{proof}

Orthogonality is the precise sense in which SP and the return-based tests answer
different questions: by construction SP is orthogonal to the exposure direction
$A$ that drives the competitors' rejections, because $A$ is invariant across
resamples and hence carries no conditional variance. The finite-sample shadow of
this is the controlled experiment of Section~\ref{sec:competitor}: in the
structural-exposure world, where a profitable strategy times its entries at
random, SP holds its nominal size at $0.050$ while Reality Check and SPA reject at
$0.660$ and $0.638$ respectively---rejecting a strategy with no timing skill
roughly two times in three---and all four tests have full power ($1.000$) against
genuine entry skill. The methods do not differ in power against a real edge; they
differ in what they treat as an edge.

\subsection{The measure family and the sensitivity range}
\label{sec:measure-family}

The preceding results condition on a single measure $\mathbb{Q}_\theta$. But the
boundary between ``structure'' (conditioned away) and ``placement'' (tested) is a
modeling choice, and different defensible choices define different nulls. We make
the conditioning measure itself part of the inferential object and ask what can be
reported once it is acknowledged to be a choice.

\begin{definition}[Admissible menu and sensitivity range]
\label{def:measure-family}
Let $\mathfrak Q=\{\mathbb{Q}_\theta:\theta\in\Theta\}$ be the family of
structure-preserving measures that all fix $(\mathcal S,\{P_t\})$ and differ only
in which additional features of the realized schedule they hold fixed---the
gap-permutation measure (conditioning on the geometry $\mathcal S$ alone, the
minimal element), the leading-gap restriction (forbidding placement in the
universal warm-up region), and the context-matched measure (restricting entries
to bars of comparable trailing volatility). Each $\mathbb{Q}_\theta$ defines a null
$H_0^\theta$ and an estimand $q_\theta=\mathbb{Q}_\theta\big(T(S)\ge
T(S_0)\big)$. The \emph{admissible menu} $\Theta_0\subseteq\Theta$ is, by
operational definition, the set of measures that condition only on features
measurable with respect to information available at entry and that preserve the
structural profile $\mathcal S$; this definition is outcome-independent---it does
not reference realized returns. Given $\Theta_0$, the data return not a scalar
``skill'' but the \emph{sensitivity range over the declared admissible menu},
$\{q_\theta:\theta\in\Theta_0\}$, summarized by the outer bounds
$[\,\min_{\theta\in\Theta_0}q_\theta,\ \max_{\theta\in\Theta_0}q_\theta\,]$.
\end{definition}

Each $q_\theta$ is point-identified given $\theta$: it is a well-defined
functional of the data once the measure is fixed. The range is a sensitivity
analysis over a finite, analyst-chosen menu, \emph{not} a sharply identified set
and \emph{not} partial identification in the Manski sense; ``the timing skill''
is not a parameter the data identify, because no $\theta$ is privileged by the
data. These are outer bounds over a finite, analyst-chosen menu, not sharp
identified sets; enlarging the menu can only widen the range. The natural
reduction of the range to a single, defensible claim is invariance across
$\Theta_0$.

\begin{definition}[Measure-invariant null and skill claim]
\label{def:invariant}
The \emph{measure-invariant null} is the intersection
$H_0^{\cap}=\bigcap_{\theta\in\Theta_0}H_0^\theta$: there is some admissible
measure under which the realized placement is exchangeable with random placement.
\emph{Measure-invariant entry-placement skill} is its complement: rejection of
$H_0^\theta$ for \emph{every} $\theta\in\Theta_0$.
\end{definition}

The claim worth making is the complement of $H_0^{\cap}$. Testing it is an
intersection--union problem: the alternative (skill under every admissible
measure) is an intersection of one-sided alternatives, so a valid test rejects
only at the worst-case, least-favorable admissible measure.

\begin{proposition}[A valid test for measure-invariant skill]
\label{prop:iut}
Let $\hat p_M^\theta$ be the Monte Carlo $p$-value of
Proposition~\ref{prop:q-validity} computed under $\mathbb{Q}_\theta$, and define
\[
p^{\mathrm{inv}}=\max_{\theta\in\Theta_0}\hat p_M^\theta .
\]
Rejecting the measure-invariant null $H_0^{\cap}$ when
$p^{\mathrm{inv}}\le\alpha$ is a level-$\alpha$ test: under $H_0^{\cap}$,
$\Pr(p^{\mathrm{inv}}\le\alpha\mid\mathcal C)\le\alpha$.
\end{proposition}

\begin{proof}
$H_0^{\cap}$ holds iff $H_0^{\theta^\ast}$ holds for at least one
$\theta^\ast\in\Theta_0$. On the event $H_0^{\theta^\ast}$,
Proposition~\ref{prop:q-validity}(i) gives
$\Pr(\hat p_M^{\theta^\ast}\le\alpha\mid\mathcal C)\le\alpha$. Since
$p^{\mathrm{inv}}=\max_\theta\hat p_M^\theta\ge\hat p_M^{\theta^\ast}$, the event
$\{p^{\mathrm{inv}}\le\alpha\}$ implies $\{\hat p_M^{\theta^\ast}\le\alpha\}$,
whence
$\Pr(p^{\mathrm{inv}}\le\alpha\mid\mathcal C)\le\Pr(\hat
p_M^{\theta^\ast}\le\alpha\mid\mathcal C)\le\alpha$.
This is the intersection--union test of \cite{LehmannRomano2005}: the maximum
$p$-value over an admissible class controls level at the least favorable element.
\end{proof}

\begin{corollary}[Power is governed by the least-favorable measure]
\label{cor:iut-power}
Under Assumption~\ref{ass:clt}, suppose the local placement alternative gives
$\mathbb{Q}_\theta$ a noncentrality $h_\theta\ge0$ as in
Proposition~\ref{prop:local-power}. Then the asymptotic power of the
intersection--union test of Proposition~\ref{prop:iut} is governed by the
least-favorable measure,
\[
\beta^{\mathrm{inv}}=\Phi\big(\min_{\theta\in\Theta_0}h_\theta-z_{1-\alpha}\big).
\]
\end{corollary}

\begin{proof}
$p^{\mathrm{inv}}\le\alpha$ requires $\hat p_M^\theta\le\alpha$ for every
$\theta$, so the rejection event is the intersection of the per-measure rejection
events. Under contiguity its probability is dominated by, and asymptotically
governed by, the measure with the smallest noncentrality, whose marginal power is
$\Phi(\min_\theta h_\theta-z_{1-\alpha})$ by Proposition~\ref{prop:local-power};
the displayed expression is the leading term.
\end{proof}

The test is therefore honestly conservative, and the conservativeness is a
\emph{price}, not a virtue: by Corollary~\ref{cor:iut-power} the power can be near
zero even if most measures have high power, because a single near-null measure
caps the noncentrality at $\min_\theta h_\theta$. ``No measure-invariant skill''
is thus a deliberately demanding, low-power bar. The price grows with the menu:
enlarging $\Theta_0$ can only lower $\min_\theta h_\theta$ and widen the
sensitivity range, so there is no free lunch in adding measures, and we do not
present the absence of an explicit $|\Theta_0|$-multiplicity correction as a
benefit---it is simply that the least-favorable element already absorbs the cost
as conservativeness. The standard's virtue is interpretive: ``significant under
every defensible measure'' is exactly the claim a referee can trust without
adjudicating which $\theta$ is correct. Two logical facts bound its
interpretation, and we state them precisely because they discipline what the
design can and cannot deliver.

\begin{remark}[Necessity is definitional; the headline is a falsification]
\label{rem:necessity}
That rejection of $H_0^\theta$ for every $\theta\in\Theta_0$ is \emph{necessary}
for a measure-invariant skill claim is definitional, not a theorem: it is the
content of Definition~\ref{def:invariant}, which defines the claim as the
conjunction. The substantive content is on the other side. The test certifies
nothing positive---unanimous rejection is \emph{not sufficient} for genuine
entry-placement skill (Proposition~\ref{prop:nonsuff})---and absence of invariance
is consistent with genuine state-conditional skill---unanimity is \emph{not
necessary} for ability (Proposition~\ref{prop:nonsuff}). The headline is therefore
a \emph{falsification}, not a certification: the design can refute a
measure-invariant timing claim but cannot establish skill.
\end{remark}

\begin{proposition}[Non-sufficiency and non-necessity]
\label{prop:nonsuff}
Let $\Theta_0$ be admissible. Then:
\emph{(i) Non-sufficiency.} Unanimous rejection is not sufficient for genuine
entry-placement skill: a confound that displaces $T(S_0)$ in the same direction
under every $\mathbb{Q}_\theta\in\mathfrak Q$ survives the intersection and is
reported as ``skill'' though it is not placement.
\emph{(ii) Non-necessity.} Unanimous rejection is not necessary for genuine
state-conditional skill: a measure $\mathbb{Q}_{\theta'}\in\Theta_0$ can absorb a
real state-conditional edge into its conditioning set, so that $H_0^{\theta'}$ is
not rejected even though placement genuinely carries information.
\end{proposition}

\begin{proof}
(i) Suppose a feature $C$ correlated with high realized return is held fixed
(hence not randomized) by every $\mathbb{Q}_\theta\in\mathfrak Q$. Then $C$
contributes identically to $T(S_0)$ and to each $T(S)$ under every measure; if $C$
also displaces $T(S_0)$ upward relative to the resampled placements---because the
realized schedule's interaction with the common conditioned feature is
favorable---then $q_\theta$ is small for all $\theta$, so $H_0^\theta$ is rejected
for all $\theta$, yet the displacement is attributable to $C$, a confound shared
by the whole class, not to placement. The intersection cannot separate a
placement effect from a common-to-all-measures confound. (ii) Take a strategy
whose edge is genuinely conditional on a market state $V$ (e.g.\
trailing-volatility regime). The context-matched measure $\mathbb{Q}_{\theta'}$
restricts resampled entries to bars matching the realized $V$-context; under
$\mathbb{Q}_{\theta'}$ the realized schedule competes only against placements
sharing its state, so the state-conditional edge is held fixed rather than tested,
and $T(S_0)$ is a central draw: $q_{\theta'}$ is not small and $H_0^{\theta'}$ is
not rejected, even though the strategy possesses genuine state-conditional timing
skill. Hence unanimity fails despite real skill.
\end{proof}

Part~(ii) is not hypothetical: Section~\ref{sec:schedule-sensitivity} exhibits
exactly such a case, in which four strategies reject under the context-matched
measure while the neutral gap-permutation measure rejects none, so a genuine
state-conditional volatility-timing edge cannot here be separated from a
matched-pool artifact. We therefore present measure invariance
(Definition~\ref{def:invariant}) as the strongest claim the design supports, not
as a characterization of skill: it is the largest claim that survives the
analyst's choice of the structure/placement partition. We adopt the
gap-permutation measure as the minimal element of $\mathfrak Q$---it adds no
contextual restriction beyond $\mathcal S$---as the primary null, without
asserting it is the unique admissible choice, and report the sensitivity range
across $\Theta_0$ rather than a single number.
